\section{Introduction}
Pre-trained models can be adapted to many downstream tasks, but
full fine-tuning updates and stores a separate copy of every model
parameter for each task. This can make adaptation expensive in optimizer
memory and checkpoint storage. Low-Rank Adaptation (LoRA) addresses this
cost by freezing the pre-trained weights and learning a small, low-rank
update for selected matrices \cite{hu2022lora}.

This report examines the parameter and deployment trade-offs of LoRA and
reviews Orthogonal LoRA (O-LoRA), which applies task-specific low-rank
updates to continual learning and penalizes overlap between task update
subspaces \cite{wang2023olora}. Our guiding question is how much LoRA
reduces trainable parameters, and what evidence is needed to assess that
efficiency alongside task quality and practical resource use.

We first review the LoRA and O-LoRA methods, then derive parameter counts
for a $1024\times1024$ matrix. We also describe a DistilRoBERTa experiment
on four-class AG News classification and discuss its reported results and
limitations.
